\chapter{总结与展望}
\label{chap:conclusion}

本章对全文研究工作进行系统总结，梳理主要贡献与核心结论，
并在此基础上指出当前工作的局限性与未来可探索的研究方向。

%% ─────────────────────────────────────────────────────────────────────────────
\section{全文总结}
\label{sec:conclusion-summary}

本文以代码修复智能体在仓库级修复任务中的高推理成本问题为核心，
围绕"成本来自哪里"与"如何系统降低成本"两个层次展开研究，
以 \texttt{mini-swe-agent} 为执行框架，以 SWE-bench Verified（500 条实例）
为评测基准，构建了完整的实验分析闭环，对六种成本优化方法与一个无增强基线进行了
系统对比。全文主要工作与贡献如下。

\subsection{建立了面向代码修复智能体的成本分析框架}
\label{subsec:contribution-framework}

第四章在形式化定义代码修复任务的基础上，提出了一套\textbf{阶段化推理成本分析框架}。
框架将 agent 的执行轨迹分解为 \texttt{SEARCH}、\texttt{READ}、\texttt{EDIT}、
\texttt{VERIFY}、\texttt{RETRY} 五个阶段，并给出了精确的 token 度量口径：

\begin{equation}
  C_{\text{total}} = \sum_{k=0}^{K-1}\bigl(c_k^{\text{in}} + c_k^{\text{out}}\bigr)
  = C_{\text{search}} + C_{\text{read}} + C_{\text{edit}} + C_{\text{verify}} + C_{\text{retry}}
\end{equation}

框架揭示了三个关键机制：
（1）对话历史累积导致输入 token 随步数超线性增长；
（2）早期读取的大文件在后续每一步都作为上下文重现，产生放大效应；
（3）失败轨迹几乎总是耗尽最大步数预算，其成本远超成功轨迹。
该框架在本文实验中通过本地透明代理网关（FastAPI 透明代理）进行了精确验证，
为后续方法设计提供了理论依据和测量工具。

\subsection{设计并实现了两条正交的成本优化路径}
\label{subsec:contribution-methods}

第五章在上述框架的指导下，围绕\textbf{上下文压缩}与\textbf{工作流压缩}
两条正交路径，设计并实现了六种优化方法：

\textbf{上下文压缩路径}（三种方法）：通过在 agent 运行前精简输入上下文，
减少搜索和阅读阶段的 token 开销。
\begin{itemize}
  \item \texttt{rag\_topk}：基于关键词检索与 BM25 打分，注入 top-4 候选文件片段
  \item \texttt{rag\_function}：在文件级检索基础上叠加函数级代码裁剪
  \item \texttt{llmlingua\_original}：使用 \texttt{deepseek-coder-1.3b} 对候选文件
        进行语义相关性重排序，替代词频统计
\end{itemize}

\textbf{工作流压缩路径}（两种方法）：通过引导 agent 遵循更高效的求解流程，
减少规划空转和无效探索带来的 token 浪费。
\begin{itemize}
  \item \texttt{skill\_abstraction}：注入六步结构化工作流脚手架，
        约束 agent 按"理解 → 定位 → 最小 patch → 验证 → 提交"路径求解
  \item \texttt{skill\_memory\_md}：注入从历史 resolved 样本中整理的
        跨任务可迁移启发式经验（\texttt{skills.md}）
\end{itemize}

\textbf{组合路径}（一种方法）：
\begin{itemize}
  \item \texttt{hybrid\_llmlingua}：串联语义重排序、函数级裁剪与 token budget 截断
\end{itemize}

所有方法均遵循统一的输入输出接口，在相同执行框架下进行公平比较。

\subsection{在 500 条实例上完成了系统实验并得出量化结论}
\label{subsec:contribution-results}

第六章基于完整 500 条 SWE-bench Verified 实例的实验结果，
得出以下核心量化结论（全部相对基线 \texttt{baseline\_raw} 计算）：

\begin{table}[htbp]
  \centering
  \caption{全文核心实验结论汇总}
  \label{tab:conclusion-summary}
  \begin{tabularx}{\textwidth}{lXrr}
    \toprule
    方法 & 主要结论 & $\Delta$Resolve & $\Delta$CPR \\
    \midrule
    \texttt{rag\_topk}          & 最高解决率（37.8\%），CPR 降低 10.0\%      & $+$3.8 pp & $-$10.0\% \\
    \texttt{skill\_memory\_md}  & 最佳综合方案，CPR 降低 10.5\%             & $+$3.2 pp & $-$10.5\% \\
    \texttt{skill\_abstraction} & 最低成功运行时 119 s，CPR 降低 9.2\%      & $+$2.4 pp & $-$9.2\%  \\
    \texttt{hybrid\_llmlingua}  & 组合压缩有效，CPR 降低 8.6\%             & $+$2.2 pp & $-$8.6\%  \\
    \texttt{llmlingua\_original}& token 削减最多（$-$5.74\%），但解决率下降  & $-$1.0 pp & $-$2.9\%  \\
    \texttt{rag\_function}      & 函数级裁剪过度，解决率和 CPR 均劣于基线   & $-$0.8 pp & $+$1.8\%  \\
    \bottomrule
  \end{tabularx}
\end{table}

三条最重要的实验发现如下：

\textbf{发现一：失败轨迹是最大的成本黑洞。}
所有失败实例几乎一律耗尽 40 步上限，平均消耗约 440--470 秒运行时；
而成功实例平均仅需 29--30 步、约 120--135 秒。
失败实例占比约 66\%，其 token 消耗约占总成本的 83\%。
这意味着降低 CPR 的最有效杠杆不是上下文压缩，而是\textbf{提升解决率或引入有效早停}。

\textbf{发现二：代码阅读（understand）阶段是 token 消耗主体（约 50\%），
代码修改（implement）仅占约 2.3\%--2.6\%。}
成本瓶颈在于"找到正确的修复位置"，而非生成 patch 本身。
因此，任何以压缩最终输出为目标的方法优化空间极为有限，
而减少冗余代码阅读的方法才更具潜力。

\textbf{发现三：不同方法优化的是不同瓶颈，具有互补性。}
\texttt{rag\_topk} 改善的是\textbf{文件定位效率}，
\texttt{llmlingua\_original} 改善的是\textbf{上下文冗余}，
\texttt{skill\_abstraction} 改善的是\textbf{工作流稳定性}，
\texttt{skill\_memory\_md} 改善的是\textbf{跨任务经验迁移}。
这一发现对未来方法组合设计具有重要指导意义。

\subsection{回答了四个核心研究问题}
\label{subsec:rq-answers}

\textbf{RQ1（成本主要来自哪些环节）：}
代码阅读（understand）阶段占总 token 消耗的 50\% 左右，是主体；
失败轨迹由于步数耗尽，其总 token 远高于成功轨迹；
搜索（search）阶段约占 10\%，验证（validate）约占 12\%，
真正的代码修改（implement）仅占 2--3\%。

\textbf{RQ2（仓库级修复为何更易引发上下文膨胀）：}
根本原因在于两点：（1）对话历史不截断导致每步输入随步数线性累积；
（2）大文件的放大效应——第 $j$ 步读取的 $F$ tokens 的文件，
在此后 $(K-j-1)$ 步中每步都会出现在输入中，总贡献为 $F \times (K-j-1)$，
远超一次阅读本身的代价。

\textbf{RQ3（已有压缩方法在代码修复中是否适用）：}
通用文本压缩方法（LLMlingua）在代码修复场景中适配性不足：
代码中看似冗余的结构性 token（缩进、类型注解、调用层次）
对语义理解具有不可替代的价值，通用压缩无法甄别其结构性价值。
函数级 RAG 因截断横向上下文而解决率下降，说明代码修复需要保留
函数簇和跨文件调用关系，而非单一函数片段。
文件级 RAG（\texttt{rag\_topk}）是迄今最适合仓库级修复的检索方案。

\textbf{RQ4（如何在保证修复能力前提下降低推理成本）：}
技能记忆注入（\texttt{skill\_memory\_md}，CPR $-$10.5\%）和
文件级 RAG（\texttt{rag\_topk}，CPR $-$10.0\%）是当前最优综合方案，
两者均同时提升解决率并降低平均 token 消耗，达到 Pareto 改善。
从机制层面，降低 CPR 的三个有效杠杆依次为：
（1）提升解决率（将失败变为成功）；
（2）减少冗余搜索与阅读（改善搜索效率）；
（3）稳定工作流（减少规划空转）。
纯粹的 token 压缩在不改善解决率的前提下，CPR 改善空间有限。


%% ─────────────────────────────────────────────────────────────────────────────
\section{工作局限性}
\label{sec:limitations}

尽管本文在 500 条 SWE-bench Verified 实例上进行了系统实验，
设计并评估了七种方法，仍存在以下几方面局限性，
对结论的适用范围有一定影响。

\subsection{单模型依赖}
\label{subsec:limit-single-model}

本文所有实验均使用 \texttt{MiniMaxAI/MiniMax-M1} 作为推理模型。
不同基础模型在代码理解能力、指令遵循能力和工作流压缩的响应方式上存在显著差异。
特别是工作流脚手架注入（\texttt{skill\_abstraction}）和经验记忆注入
（\texttt{skill\_memory\_md}）的有效性，在很大程度上取决于模型对结构化提示的
理解与遵循能力。在更强或更弱的基础模型上，本文方法的效果可能存在一定偏差。
未来工作需要在 GPT-4o、Claude 3.5 Sonnet、DeepSeek-V3 等代表性模型上
进行跨模型验证，以评估方法的模型无关性。

\subsection{评估基准的代表性局限}
\label{subsec:limit-benchmark}

SWE-bench Verified 500 条实例主要覆盖 \texttt{astropy}、
\texttt{django}、\texttt{sympy}、\texttt{matplotlib}、\texttt{scikit-learn}
等 Python 开源仓库，且以函数级 bug 修复为主。
这一分布可能无法完全代表工业界真实软件仓库中的修复任务，
例如：跨语言代码库、遗留代码系统、与数据库或外部服务深度耦合的业务逻辑 bug 等。
此外，SWE-bench 的任务难度分布并非均匀：约 40\% 的实例（203 条）对全部方法均失败，
这些"硬实例"的存在使得 resolve rate 的绝对数值对方法优劣的区分度降低，
也限制了 CPR 指标在极端情况下的解释能力。

\subsection{BM25 检索的关键词依赖性}
\label{subsec:limit-retrieval}

本文的文件级 RAG 检索基于关键词提取与 BM25 打分，
其效果高度依赖问题描述中是否包含可识别的函数名、类名或模块路径等显式线索。
对于问题描述较为模糊、以自然语言描述现象而非直接引用代码符号的实例，
BM25 检索可能无法召回正确文件，从而导致 \texttt{rag\_topk} 退化为噪声注入。
更高质量的语义检索（如基于 code embedding 的向量检索）可能在这类实例上
取得更好效果，但本文未进行系统比较。

\subsection{skill memory 的手工构建成本}
\label{subsec:limit-skill}

\texttt{skill\_memory\_md} 依赖人工整理的 \texttt{skills.md} 文件，
其内容是基于部分已 resolved 样本总结的启发式规则。
这一构建过程存在以下局限：
（1）经验规则覆盖的模式有限，对规则未覆盖的新类型 bug 帮助有限；
（2）构建过程存在数据泄露风险（若规则来源样本与测试样本有重叠），
本文实验在设计上保证了训练集与测试集的隔离，但工业应用中需要更严格的审计；
（3）当基础模型更换时，已有 skill memory 的有效性可能需要重新评估。
自动化的 skill memory 构建与持续更新机制，是未来将该方法推向实用的关键挑战。

\subsection{缺乏早停机制}
\label{subsec:limit-earlystop}

本文实验固定最大步数为 $K=40$，所有失败实例均消耗全部步数预算。
然而，实验结果表明失败轨迹的平均 token 消耗是成功轨迹的 2--3 倍，
而失败与成功的分化在轨迹早期往往已有信号可循（例如反复读取同一文件、
patch 格式错误循环、搜索无收敛等）。
本文尚未研究基于轨迹特征的早停策略，这是当前成本瓶颈中
"失败成本过高"问题最直接的解决方向，但也需要在不误判成功轨迹的前提下
设计合适的终止条件，具有一定的实现复杂度。


%% ─────────────────────────────────────────────────────────────────────────────
\section{未来工作展望}
\label{sec:future-work}

基于本文的实验发现与局限性分析，以下几个方向具有较高的研究价值与工程落地潜力。

\subsection{轨迹级早停与无进展检测}
\label{subsec:future-earlystop}

当前成本的最大黑洞是失败轨迹的步数耗尽问题。
未来工作可以从以下两个角度设计轨迹级成本控制机制：

\textbf{（1）无进展检测（stagnation detection）。}
通过分析连续步骤的操作类型、文件读取重复率和 patch 尝试的差异性，
识别轨迹陷入循环的模式。当检测到无进展信号时触发策略切换（如扩展搜索范围、
换用不同提示策略）或提前终止。

\textbf{（2）置信度估计（confidence estimation）。}
在每一步生成后估计当前状态下任务成功的概率，
当概率低于阈值时优雅退出并返回当前最优 patch（若有），
而非继续消耗预算。这需要训练一个轻量级的 trajectory-level 状态分类器。

初步估计，若将失败轨迹的平均步数从 40 步降至 25 步，
在 resolve rate 不变的前提下，全局平均 CPR 可改善约 20--30\%，
优于本文所有已实现方法的改善幅度。

\subsection{自动化 skill memory 构建与更新}
\label{subsec:future-skill-memory}

当前 \texttt{skill\_memory\_md} 依赖人工整理的启发式规则，
限制了其规模化应用能力。未来可以探索以下自动化路径：

\textbf{（1）从成功轨迹自动提炼经验。}
对已 resolved 实例的 \texttt{trajectory.json} 进行结构化分析，
自动提取关键决策模式（如"在 X 类问题上应优先搜索 Y 类文件"），
并用语言模型将其归纳为可复用的文字规则，动态更新 \texttt{skills.md}。

\textbf{（2）基于强化学习的 skill 演化。}
将技能规则视为策略参数，通过在新实例上的执行反馈（resolved/unresolved）
评估每条规则的贡献，并用梯度信号或进化算法迭代优化规则集合。

\textbf{（3）分仓库、分 bug 类型的专属 skill memory。}
不同仓库的代码结构和 bug 分布存在显著差异，
为 \texttt{django}、\texttt{astropy} 等高频仓库单独维护 skill memory
可能比通用 memory 更加有效。

\subsection{语义向量检索替代 BM25}
\label{subsec:future-embedding}

当前 \texttt{rag\_topk} 的文件检索基于 BM25 关键词匹配，
对问题描述依赖显式代码标识符的实例效果好，
但对语义相关性不依赖关键词共现的问题描述存在召回率下降的问题。
未来可以引入基于 code embedding 的向量检索，例如：

\begin{itemize}
  \item 使用 \texttt{CodeBERT}、\texttt{UnixCoder} 或
        \texttt{text-embedding-3-large} 对仓库文件进行离线向量化
  \item 在 agent 运行前，用问题描述的 embedding 查询 Top-K 文件
  \item 将向量检索与 BM25 结果进行融合重排（RRF，Reciprocal Rank Fusion）
\end{itemize}

语义检索尤其对"问题描述描述的是现象而非代码符号"的实例更有帮助，
有望进一步提升 \texttt{rag\_topk} 在当前 37.8\% 基础上的解决率。

\subsection{分层组合方法设计}
\label{subsec:future-pipeline}

实验结果表明，不同方法优化的是不同瓶颈，具有互补性：
\texttt{rag\_topk} 改善文件定位效率，
\texttt{skill\_memory\_md} 改善决策与执行质量。
一个自然的后续方向是将这两条路径系统地组合为三层 pipeline：

\begin{enumerate}
  \item \textbf{检索层}：利用向量检索或 BM25 缩小候选文件集合（对标 \texttt{rag\_topk}）
  \item \textbf{决策层}：注入结构化工作流与可迁移经验
        （对标 \texttt{skill\_abstraction} + \texttt{skill\_memory\_md}）
  \item \textbf{止损层}：加入基于轨迹特征的无进展检测与早停机制
\end{enumerate}

三层方法分别对应"定位"、"收敛"、"止损"三个瓶颈，
理论上可将 CPR 改善叠加到 20\% 以上，
同时避免单一方法过拟合某类实例的风险。

\subsection{跨模型泛化性验证}
\label{subsec:future-model-transfer}

本文所有实验基于 \texttt{MiniMax-M1} 模型。
为验证所提方法的普适性，未来工作应在以下模型上进行跨模型对比：
\begin{itemize}
  \item 开源模型：\texttt{DeepSeek-V3}、\texttt{Qwen2.5-Coder-32B}、\texttt{Llama-3.3-70B}
  \item 商业模型：\texttt{GPT-4o}、\texttt{Claude-3.5-Sonnet}
\end{itemize}

重点验证的问题包括：
（1）工作流脚手架（\texttt{skill\_abstraction}）对不同指令遵循能力模型的效果差异；
（2）BM25 检索注入对模型上下文利用率的影响是否模型无关；
（3）skill memory 在不同模型上是否需要重新调优。

\subsection{面向工业代码库的适配性研究}
\label{subsec:future-industrial}

SWE-bench 任务以 Python 开源库为主，
与工业界的真实软件修复场景存在一定差异：
工业代码库通常规模更大、多语言混合、依赖关系更复杂、
问题描述的质量与格式也更加多样。
未来工作可以从以下方面推进工业适配性研究：
\begin{itemize}
  \item 在 SWE-bench-Java 或自建多语言基准上验证方法的跨语言适应性
  \item 研究在大型单体仓库（百万行代码级别）中的检索效率与覆盖率
  \item 针对问题描述质量差（无代码标识符、描述模糊）的实例设计专项增强策略
  \item 结合持续集成系统（CI/CD）设计面向工业部署的成本预算控制机制
\end{itemize}


%% ─────────────────────────────────────────────────────────────────────────────
\section{本章小结}
\label{sec:chap07-summary}

本章对全文研究工作进行了系统总结与展望。
在研究贡献方面，本文围绕代码修复智能体的高推理成本问题，
建立了阶段化成本分析框架，提出并实现了上下文压缩与工作流压缩两条正交优化路径，
在 500 条 SWE-bench Verified 实例上完成了覆盖七种方法的系统实验，
得出了若干具有实践指导意义的量化结论。

核心结论可以用一句话概括：
\textbf{代码修复智能体的成本瓶颈不在于生成补丁，而在于"搜索正确修复位置"的过程；
改善这一过程的最有效手段是提升解决率与减少失败轨迹的无效消耗，
而非简单地压缩输入 token。}

在局限性方面，本文工作受限于单模型依赖、评估基准分布偏差、
手工 skill memory 的规模瓶颈与缺乏早停机制。
这些局限性同时也指向了最有价值的未来研究方向：
轨迹级早停、自动化 skill 构建、语义向量检索、分层 pipeline 设计
以及跨模型泛化性验证。

本文的工作表明，代码修复智能体的工程部署可行性不仅取决于模型能力的提升，
同样依赖于对执行过程的精细化管控与成本意识的系统化设计。
随着大语言模型在软件工程自动化中的应用不断深入，
对推理成本的系统研究将成为该领域不可回避的重要课题。
